% Generated draft: requires mathematical and source review.
\documentclass[11pt]{article}
\usepackage[margin=25mm]{geometry}
\usepackage{fontspec}
\setmainfont{Latin Modern Roman}
\usepackage{amsmath,amssymb,mathtools,unicode-math}
\setmathfont{Latin Modern Math}
\usepackage{enumitem,needspace,xurl,hyperref}
\setlength{\parindent}{0pt}
\setlength{\parskip}{5pt}
\setlength{\emergencystretch}{3em}
\newcommand{\N}{\mathbb{N}}
\newcommand{\Z}{\mathbb{Z}}
\newcommand{\Q}{\mathbb{Q}}
\newcommand{\R}{\mathbb{R}}
\newcommand{\C}{\mathbb{C}}
\newcommand{\F}{\mathbb{F}}
\newcommand{\A}{\mathbb{A}}
\newcommand{\PP}{\mathbb{P}}
\newcommand{\E}{\mathbb{E}}
\newcommand{\eps}{\varepsilon}
\newcommand{\dd}{\,\mathrm{d}}
\begin{document}
\section*{UnsolvedMath ID: 701007}
\textbf{Source problem code:} catalogue entry
\subsection{Definitions and mathematical statement}
Let \(G\) be a torsion-free linear group: thus \(g^m=1\), for \(g\in G\) and an integer \(m\geq 1\), implies \(g=1\), and \(G\) embeds in \(\operatorname{GL}_d(K)\) for some field \(K\) and some positive integer \(d\). The source imposes no restriction on the characteristic of \(K\). Let \(X=\{x_1,\ldots,x_n\}\) be a finite set of variables, let \(F(X)\) be the free group on \(X\), and put \(G[X]=G\ast F(X)\).

For any subset \(U\subseteq G[X]\), define its solution set over \(G\) by
\[
V_G(U)=\{\mathbf g\in G^n: u(\mathbf g)=1\text{ for every }u\in U\},
\]
where evaluation fixes the coefficient subgroup \(G\) and sends each \(x_i\) to \(g_i\). Two systems are equivalent over \(G\) if they have the same solution set. Define
\[
\operatorname{Rad}_G(U)=\{w\in G[X]:w(\mathbf g)=1\text{ for every }\mathbf g\in V_G(U)\}.
\]
A subgroup \(H\leq G[X]\) is isolated if, for every \(w\in G[X]\) and every integer \(m\geq1\), the inclusion \(w^m\in H\) implies \(w\in H\). Define
\[
\sqrt{U}=\bigcap\{N\trianglelefteq G[X]: U\subseteq N\text{ and }N\text{ is isolated}\}.
\]
This is the smallest normal isolated subgroup containing \(U\).

The source does not define an irreducible system. For this draft, use the conventional interpretation that \(V_G(S)\) is nonempty and irreducible in the Zariski topology on \(G^n\), namely the topology whose closed sets are generated by the solution sets \(V_G(U)\) under finite unions and arbitrary intersections. Irreducibility means that \(V_G(S)\) cannot be written as the union of two proper relatively closed subsets. This interpretation requires confirmation against the source's intended convention.

\textbf{Requested assertion (catalogue entry).} For every torsion-free linear group \(G\) and every subset \(S\subseteq G[X]\) defining an irreducible system, there exists a finite subset \(T\subseteq G[X]\) such that
\[
V_G(S)=V_G(T)
\qquad\text{and}\qquad
\operatorname{Rad}_G(T)=\sqrt{T}.
\]
The system \(S\) is not assumed finite. The supplied source reports that the assertion holds when \(G\) is a free group, citing O. Kharlampovich and A. Myasnikov (1998). It also records explicitly unverified AI claims that the general assertion is false. Accordingly, the displayed assertion is the original proposed target, not an established theorem or an independently verified assessment of its current status.
\subsection{Short English statement}
The problem asks whether every irreducible system of equations over a torsion-free linear group has an equivalent finite system whose equations imply exactly the words in its smallest normal isolated subgroup. The intended meaning of irreducibility needs confirmation. The supplied registry labels the problem open while recording unverified claims of counterexamples; this draft preserves both qualifications.
\subsection{Sources}
\begin{itemize}
\item Local UnsolvedMath registry, supplied complete record: permanent ID \(701007\), source code \texttt{catalogue entry}, title \emph{Kourovka Notebook Problem 14.22}. Upstream revision: \texttt{d371ee917ac9647fbb06cd6fc529deec6f3dda82}.
\item G. Baumslag, A. G. Myasnikov, and V. N. Remeslennikov, Problem 14.22, \emph{The Kourovka Notebook (Unsolved Problems in Group Theory)}, 14th issue (1999); also identified in the supplied record with the 21st issue edition (2026), \texttt{arXiv:1401.0300v48}. The registry identifies E. I. Khukhro and V. D. Mazurov as editors of the collection.
\item O. Kharlampovich and A. Myasnikov, \emph{J. Algebra} \textbf{200} (1998), 472--570. The supplied problem statement cites this reference for the free-group case; no article title is supplied, and the claim has not been independently checked here.
\item Achyuth Jayadevan, \emph{Preprint of 9 September 2026}, \url{https://zenodo.org/records/22722075}. The supplied discussion labels this an unverified AI solution claim supporting the assertion that the problem is false in general.
\item M. Kinyon and J. Urban, \emph{Preprint of 17 September 2026}, \url{https://kourovkanotebookorg.wordpress.com/wp-content/uploads/2026/09/kourovka-mathematics.pdf}. The supplied discussion likewise labels this an unverified AI solution claim.
\item Supplied status review dated \texttt{2026-10-16}: canonical status \texttt{open}, evidence category \texttt{unrefereed\_claim}. It reports a \(\#\) marker for an unverified AI solution claim in the October 2026 21st issue, \texttt{arXiv:1401.0300v48}. This is supplied dated curation, not an independently verified literature assessment. The status-review note truncates its second URL; the complete URL above is supplied separately in the background field.
\end{itemize}
\end{document}
